\documentclass[11pt]{article}
\usepackage[margin=2.5cm]{geometry}
\usepackage{amsmath,amssymb,amsthm}
\usepackage{mathtools}
\usepackage{url}
\usepackage{hyperref}

\newtheorem{theorem}{Theorem}
\newtheorem{lemma}[theorem]{Lemma}
\newtheorem{proposition}[theorem]{Proposition}
\newtheorem{corollary}[theorem]{Corollary}
\newtheorem{definition}[theorem]{Definition}
\newtheorem{remark}[theorem]{Remark}
\newtheorem{problem}[theorem]{Problem}

\newcommand{\R}{\mathbb{R}}
\newcommand{\Q}{\mathbb{Q}}
\newcommand{\N}{\mathbb{N}}
\newcommand{\evalu}{\widehat{p}^{+}}
\newcommand{\evalv}{\widehat{p}^{-}}
\newcommand{\absp}{\lvert p \rvert}
\newcommand{\abs}[1]{\lvert #1 \rvert}
\newcommand{\norm}[1]{\lVert #1 \rVert}

\title{Constant-Rate Dual-Railing of Bounded Polynomial GPACs:\\
       UCNC 2025 Problem~1, General $n$-Dimensional Case}
\date{\today}

\begin{document}
\maketitle

\begin{abstract}
The selective dual-railing construction of Haisler, Huang, Migunov,
Mohammed, and Provence (UCNC 2025) converts a polynomial vector field
$\dot y_i = p_i(y)$ into a chemical reaction network on
$2n$ non-negative species $(u_1, v_1, \dots, u_n, v_n)$ with
$y_i = u_i - v_i$.  A shared annihilation term $-k\,u_i v_i$ prevents
each rail-pair from drifting apart.  UCNC~2025 Problem~1 asks: given a
bounded polynomial GPAC trajectory, does there exist a single constant
rate $k$ (independent of the bound) for which the dual-railed system
remains bounded on $[0, \infty)$?  We answer affirmatively for the
general $n$-dimensional case starting from the dual-rail split of the
GPAC initial condition.  Concretely: for every polynomial vector field
$p\colon \R^n \to \R^n$ admitting a bounded solution $y_{\mathrm{Sol}}$
with $\norm{y_{\mathrm{Sol}}(t)}_\infty \le \beta$, the rate
\[
  k_\star \;:=\; \max_i \frac{2 M_i(\beta)}{3\beta^2},
  \qquad
  M_i(\beta) \;:=\; \max_{w \in [0, 2\beta]^{2n}}
                    \bigl(\evalu_i(w) + \evalv_i(w)\bigr),
\]
suffices: any rational $k > k_\star$ keeps the dual-rail trajectory
inside the box $\{(u,v) \ge 0 : u_i + v_i \le 2\beta\}$, and the
difference $u_i - v_i$ recovers $y_{\mathrm{Sol}, i}$ exactly by ODE
uniqueness.  The construction has been formalized in Lean~4 with no
unproved goals and no axioms beyond Lean's standard logical core
(propositional extensionality, the axiom of choice, and the soundness
of quotient types).
\end{abstract}

%---------------------------------------------------------------------
\section{Setup}
%---------------------------------------------------------------------

\paragraph{Polynomial GPAC.}
Fix $n \ge 1$ and a tuple
$p = (p_1, \dots, p_n) \in \Q[X_1, \dots, X_n]^n$.  We write
$\widehat{p}_i^{+}, \widehat{p}_i^{-} \in \Q[U_1, V_1, \dots, U_n, V_n]$
for the positive and negative parts of the dual-rail substitution
$y_j \mapsto U_j - V_j$ in $p_i$: every monomial of $p_i$ with positive
coefficient lands in $\widehat{p}_i^{+}$ after binomial expansion, every
monomial with negative coefficient lands (with sign flipped) in
$\widehat{p}_i^{-}$.  Both have non-negative coefficients.  The
fundamental identity
\begin{equation}\label{eq:diff-id}
  \widehat{p}_i^{+}(u, v) - \widehat{p}_i^{-}(u, v)
    \;=\; p_i(u - v),
\end{equation}
holds in $\Q[U, V]$ and is the algebraic engine of the construction.

\paragraph{Constant-rate dual-rail system.}
For a positive rational $k$, the \emph{constant-rate dual-rail
GPAC associated to $(p, k)$} is the polynomial vector field on
$\R^{2n}$:
\begin{equation}\label{eq:dr-system}
  \begin{cases}
    \dot u_i \;=\; \widehat{p}_i^{+}(u, v) - k\, u_i v_i, \\
    \dot v_i \;=\; \widehat{p}_i^{-}(u, v) - k\, u_i v_i,
  \end{cases}
  \qquad
  i = 1, \dots, n.
\end{equation}
With non-negative coefficients and a single per-pair degradation, this
is mass-action implementable as a CRN.

\paragraph{Bounded GPAC.}
A \emph{bounded GPAC datum} is a tuple
$(p, y_0, y_{\mathrm{Sol}}, \beta)$ with $\beta > 0$,
$y_0 \in \Q^n$, and $y_{\mathrm{Sol}}\colon \R \to \R^n$ such that
$y_{\mathrm{Sol}}(0) = y_0$ (cast to reals),
$\dot y_{\mathrm{Sol}, i}(t) = p_i(y_{\mathrm{Sol}}(t))$ for $t \ge 0$,
and $\abs{y_{\mathrm{Sol}, i}(t)} \le \beta$ for $t \ge 0$ and every
$i$.  In particular $\abs{(y_0)_i} \le \beta$.

\paragraph{Dual-rail split of the initial condition.}
Define $u_0, v_0 \in \Q_{\ge 0}^n$ by
\[
  (u_0)_i := \max((y_0)_i, 0), \qquad
  (v_0)_i := \max(-(y_0)_i, 0),
\]
so $(u_0)_i - (v_0)_i = (y_0)_i$ and
$(u_0)_i + (v_0)_i = \abs{(y_0)_i} \le \beta$.

\paragraph{Box.}
The forward-invariant set we will exhibit is
\[
  B_\beta \;:=\; \bigl\{ w \in \R^{2n}_{\ge 0}\;:\;
                  w_{2i} + w_{2i+1} \le 2\beta\,\,\forall i \bigr\}.
\]
We index the $2n$ dual-rail coordinates by $K \in \{0, \dots, 2n-1\}$
with $u_i = w_{2i}$ and $v_i = w_{2i+1}$.

%---------------------------------------------------------------------
\section{Main theorem}
%---------------------------------------------------------------------

\begin{theorem}[UCNC~2025 Problem~1, general $n$]\label{thm:main}
Let $(p, y_0, y_{\mathrm{Sol}}, \beta)$ be a bounded polynomial GPAC
datum and let
\[
  k_\star \;:=\; \max_{i \in \{1,\dots,n\}}\,
                 \frac{2 M_i(\beta)}{3 \beta^2},
  \qquad
  M_i(\beta) \;:=\; \sup_{w \in [0, 2\beta]^{2n}}
                    \bigl(\widehat{p}_i^{+}(w) + \widehat{p}_i^{-}(w)\bigr).
\]
For every rational $k$ with $k > k_\star$, there exists a global
solution $\widehat y\colon [0, \infty) \to \R^{2n}$ of~\eqref{eq:dr-system}
with $\widehat y(0) = (u_0, v_0)$ such that, for every $t \ge 0$,
\begin{enumerate}
\item[\textup{(i)}]\textbf{Non-negativity.}
$\widehat y(t)_K \ge 0$ for every $K$.
\item[\textup{(ii)}]\textbf{$\sigma$-bound.}
$\widehat y(t)_{2i} + \widehat y(t)_{2i+1} \le 2\beta$ for every $i$.
\item[\textup{(iii)}]\textbf{Dual-rail identity.}
$\widehat y(t)_{2i} - \widehat y(t)_{2i+1} = y_{\mathrm{Sol}, i}(t)$
for every $i$.
\end{enumerate}
\end{theorem}

The bound (ii) combined with (i) yields the box invariance
$\widehat y(t) \in B_\beta$, and (iii) closes the dual-rail
implementation: $\widehat y_{2i} - \widehat y_{2i+1}$ tracks the
original GPAC trajectory exactly, so the CRN is a faithful simulator of
$y_{\mathrm{Sol}}$.

The proof decomposes into four sub-lemmas, each formally stated in the
Lean development.  Sections~\ref{sec:identities}--\ref{sec:identity}
walk through them.

%---------------------------------------------------------------------
\section{Algebraic identities}\label{sec:identities}
%---------------------------------------------------------------------

\paragraph{Drift difference.}
Subtracting the two lines of~\eqref{eq:dr-system} and using
\eqref{eq:diff-id}:
\begin{equation}\label{eq:drift-diff}
  \dot u_i - \dot v_i
    \;=\; \widehat{p}_i^{+}(u,v) - \widehat{p}_i^{-}(u,v)
    \;=\; p_i(u - v).
\end{equation}
The annihilation term cancels.  Hence the difference $z_i := u_i - v_i$
satisfies the original GPAC.

\paragraph{Drift sum on the box.}
Adding the two lines of~\eqref{eq:dr-system}:
\[
  \dot u_i + \dot v_i
    \;=\; \widehat{p}_i^{+}(u,v) + \widehat{p}_i^{-}(u,v)
          - 2 k\, u_i v_i.
\]
Since both $\widehat{p}_i^{+}$ and $\widehat{p}_i^{-}$ have
non-negative coefficients, the absolute polynomial
$\absp_i := \widehat{p}_i^{+} + \widehat{p}_i^{-}$ is monotonic in each
coordinate over $\R^{2n}_{\ge 0}$.  On the box $[0, 2\beta]^{2n}$,
$\absp_i$ is therefore bounded by $M_i(\beta)$ (its supremum on the
box).  Hence on $B_\beta$ we have the upper bound
\begin{equation}\label{eq:drift-sum}
  \dot u_i + \dot v_i
    \;\le\; M_i(\beta) - 2 k\, u_i v_i.
\end{equation}

\paragraph{Face inequality.}
A short algebraic computation shows that on the upper face $u_i + v_i =
2\beta$ subject to $\abs{u_i - v_i} \le \beta$, the product $u_i v_i$ is
bounded below:
\begin{equation}\label{eq:uv-face}
  4 u_i v_i \;=\; (u_i + v_i)^2 - (u_i - v_i)^2
              \;\ge\; (2\beta)^2 - \beta^2
              \;=\; 3\beta^2.
\end{equation}

\paragraph{Strict inward-pointing at the upper face.}
Combining \eqref{eq:drift-sum} and \eqref{eq:uv-face}, on the face
$u_i + v_i = 2\beta$ inside $B_\beta$ with the additional dual-rail
constraint $\abs{u_i - v_i} \le \beta$,
\[
  \dot u_i + \dot v_i
    \;\le\; M_i(\beta) - 2k \cdot \tfrac{3}{4}\beta^2
    \;=\; M_i(\beta) - \tfrac{3}{2}k\beta^2.
\]
This is strictly negative whenever
$k > 2 M_i(\beta) / (3 \beta^2)$, which is exactly the
threshold condition $k > k_\star$.  Hence:
\begin{lemma}[upper-face strict negativity]\label{lem:strict-neg}
If $k > k_\star$, $\beta > 0$, $w \in B_\beta$ with
$w_{2i} + w_{2i+1} = 2\beta$ and $\abs{w_{2i} - w_{2i+1}} \le \beta$,
then the $\sigma$-drift $\dot u_i + \dot v_i$ at $w$ is strictly less
than zero.
\end{lemma}

%---------------------------------------------------------------------
\section{Picard existence}\label{sec:picard}
%---------------------------------------------------------------------

The right-hand side of~\eqref{eq:dr-system} is polynomial, hence locally
Lipschitz on $\R^{2n}$.  By the standard Picard--Lindel\"of argument,
combined with the a~priori bound $\widehat y(t) \in B_\beta$ established
in the next section, there exists a global solution
$\widehat y\colon [0, \infty) \to \R^{2n}$ of~\eqref{eq:dr-system} with
$\widehat y(0) = (u_0, v_0)$.

%---------------------------------------------------------------------
\section{The invariant: nonneg + diff + $\sigma$-bound}\label{sec:invariant}
%---------------------------------------------------------------------

We prove the three-pronged invariant on every interval $[0, T)$ on
which the candidate solution exists.  The three prongs --- non-negativity,
difference bound, and $\sigma$-bound --- are proved in this order and
assembled in Section~\ref{sec:assembly}.

%---------------------------------------------------------------------
\subsection{Componentwise non-negativity}
%---------------------------------------------------------------------

\begin{lemma}\label{lem:nonneg}
Let $\widehat y\colon \R \to \R^{2n}$ satisfy~\eqref{eq:dr-system} on
$[0, T)$ with non-negative initial condition.  Then $\widehat y(t)_K
\ge 0$ for all $t \in [0, T)$ and all~$K$.
\end{lemma}

\begin{proof}
The system~\eqref{eq:dr-system} admits a CRN decomposition: each
component $\dot w_K = (\text{production}) - k\, w_K w_{K'}$ where the
production term is $\widehat{p}_i^{+}$ or $\widehat{p}_i^{-}$ (both
non-negative on $\R^{2n}_{\ge 0}$), and the degradation $-k w_K w_{K'}$
is the product of $w_K$ with another component.  At any face $w_K = 0$
the production is non-negative and the degradation vanishes, so the
face is inward-pointing.  The standard local non-negativity argument
for CRNs then yields $\widehat y(t) \ge 0$ for all $t$ in the lifetime
of the solution.
\end{proof}

%---------------------------------------------------------------------
\subsection{Difference bound by ODE uniqueness}
%---------------------------------------------------------------------

\begin{lemma}\label{lem:diff}
Under the hypotheses of Theorem~\ref{thm:main}, on $[0, T)$,
\[
  \abs{\widehat y(t)_{2i} - \widehat y(t)_{2i+1}}
    \;\le\; \beta.
\]
\end{lemma}

\begin{proof}
Let $z_i(t) := \widehat y(t)_{2i} - \widehat y(t)_{2i+1}$.  By
\eqref{eq:drift-diff}, $\dot z = p(z)$ on $[0, T)$.  At $t = 0$:
\[
  z_i(0)
    = (u_0)_i - (v_0)_i
    = \max((y_0)_i, 0) - \max(-(y_0)_i, 0)
    = (y_0)_i
    = y_{\mathrm{Sol}, i}(0).
\]
Both $z$ and $y_{\mathrm{Sol}}$ satisfy the original polynomial GPAC
$\dot w = p(w)$ with the same initial condition.  On any compact
sub-interval $[0, t] \subset [0, T)$, both solutions are bounded
(continuity on a compact interval) and $p$ is locally Lipschitz, so
the standard ODE uniqueness theorem on a compact interval gives
$z = y_{\mathrm{Sol}}$ on $[0, t]$.  Letting $t \uparrow T$:
$z(t) = y_{\mathrm{Sol}}(t)$ on $[0, T)$.  By the bound assumption,
$\abs{y_{\mathrm{Sol}, i}(t)} \le \beta$.
\end{proof}

\begin{remark}
This step does \emph{not} use the $\sigma$-bound and so can be carried
out independently.  Critically, this is what breaks a potential
circular dependency: the Nagumo $\sigma$-barrier needs the diff bound
inside the strict-negativity lemma~\ref{lem:strict-neg}, but the diff
bound itself uses only ODE uniqueness on compact subintervals (with
sup-norm bounds extracted by compactness, not by $B_\beta$).
\end{remark}

%---------------------------------------------------------------------
\subsection{The joint $\sigma$-Nagumo barrier}\label{sec:sigma}
%---------------------------------------------------------------------

This is the analytic core.  Define
$\sigma_i(t) := \widehat y(t)_{2i} + \widehat y(t)_{2i+1}$.

\begin{lemma}\label{lem:sigma}
Under the hypotheses of Theorem~\ref{thm:main} and assuming
$k > k_\star$, on $[0, T)$,
\[
  \sigma_i(t) \;\le\; 2\beta \qquad \forall i.
\]
\end{lemma}

\begin{proof}
By Lemmas~\ref{lem:nonneg} and~\ref{lem:diff}, $\widehat y_K \ge 0$ and
$\abs{z_i} \le \beta$ on $[0, T)$.  Suppose for contradiction there
exist $t' \in [0, T)$ and $i_0$ with $\sigma_{i_0}(t') > 2\beta$.
Define
\[
  S \;:=\; \bigl\{ \tau \in [0, t']\,:\,
            \sigma_i(\tau) \le 2\beta\,\,\forall i \bigr\}.
\]

\smallskip
\noindent
\textbf{$0 \in S$.}  At $t = 0$, $\sigma_i(0) = (u_0)_i + (v_0)_i =
\abs{(y_0)_i} \le \beta < 2\beta$.

\smallskip
\noindent
\textbf{Define $\tau_\star := \sup S \in [0, t']$.}  By sequential
continuity of each $\sigma_i$, $\tau_\star \in S$, hence
$\sigma_i(\tau_\star) \le 2\beta$ for every $i$.

\smallskip
\noindent
\textbf{$\tau_\star < t'$.}  By assumption $\sigma_{i_0}(t') > 2\beta$,
so $t' \notin S$ and $\tau_\star < t'$.

\smallskip
\noindent
\textbf{Box property at $\tau_\star$.}  Since $\widehat y_K(\tau_\star)
\ge 0$ and $\sigma_i(\tau_\star) \le 2\beta$, each component
$\widehat y_K(\tau_\star)$ is bounded above by its $\sigma$-partner's
sum minus the other component, which is at most $\sigma_i(\tau_\star)
\le 2\beta$.  Hence $\widehat y(\tau_\star) \in B_\beta$.

\smallskip
\noindent
\textbf{Per-coordinate buffers.}  For each $i$:
\begin{itemize}
\item If $\sigma_i(\tau_\star) < 2\beta$: by continuity of $\sigma_i$,
there exists $\delta_i > 0$ such that
$\sigma_i(\tau) \le 2\beta$ on $[\tau_\star, \tau_\star + \delta_i]$.
\item If $\sigma_i(\tau_\star) = 2\beta$: by Lemma~\ref{lem:strict-neg}
(applied with $w = \widehat y(\tau_\star)$, using the box property,
the face condition, and Lemma~\ref{lem:diff}),
$\dot \sigma_i(\tau_\star) < 0$.  A function with strictly negative
derivative at a point lies below its value on a right neighborhood:
explicitly, by the definition of the derivative, for some
$\delta_i > 0$ and all $\tau \in [\tau_\star, \tau_\star + \delta_i]$,
\[
  \sigma_i(\tau) \;\le\; \sigma_i(\tau_\star) + \tfrac{1}{2}\,
                          \dot\sigma_i(\tau_\star)\,(\tau - \tau_\star)
                 \;\le\; \sigma_i(\tau_\star) \;=\; 2\beta.
\]
\end{itemize}

\smallskip
\noindent
\textbf{Uniform buffer.}  Take
$\delta := \min\bigl(\min_i \delta_i,\, t' - \tau_\star\bigr) > 0$.
The minimum over $i$ is well-defined and positive because
$\{1, \dots, n\}$ is finite and every $\delta_i > 0$.  Then on
$[\tau_\star, \tau_\star + \delta]$, all $\sigma_i \le 2\beta$, hence
$\tau_\star + \delta \in S$.

\smallskip
\noindent
\textbf{Contradiction.}  But $\tau_\star + \delta > \tau_\star = \sup S$,
contradicting the definition of supremum.\qedhere
\end{proof}

\begin{remark}[Why the barrier is joint]
A naive per-coordinate barrier --- defining
$\tau_i := \sup\{\tau : \sigma_i(\tau) \le 2\beta\}$ separately for
each $i$ --- fails because Lemma~\ref{lem:strict-neg}'s box hypothesis
$w \in B_\beta$ requires \emph{all} coordinates to satisfy
$w_{2j} + w_{2j+1} \le 2\beta$, not just the $i$-th pair.  Choosing
$i_\star$ as the last to cross does not help, because by that time
other $\sigma_j$ may already have escaped.  The joint sup
$\tau_\star := \sup\{\tau : \forall i,\,\sigma_i(\tau) \le 2\beta\}$
is the correct barrier: at $\tau_\star$ all $\sigma_i$ are
simultaneously $\le 2\beta$, hence $w(\tau_\star) \in B_\beta$, hence
the strict inward-pointing condition applies to every $i$ saturating
the upper face.
\end{remark}

%---------------------------------------------------------------------
\section{Closing the dual-rail identity}\label{sec:identity}
%---------------------------------------------------------------------

\begin{lemma}\label{lem:identity}
Under the hypotheses of Theorem~\ref{thm:main}, on $[0, \infty)$,
\[
  \widehat y(t)_{2i} - \widehat y(t)_{2i+1}
    \;=\; y_{\mathrm{Sol}, i}(t).
\]
\end{lemma}

This is what we already proved as Lemma~\ref{lem:diff}, but now over
the global interval $[0, \infty)$ (not just on a finite $[0, T)$).

\begin{proof}
The argument of Lemma~\ref{lem:diff} applies verbatim with $T =
\infty$, since the global Picard solution $\widehat y$ exists on
$[0, \infty)$ and ODE uniqueness on every compact $[0, t]$ pins the
difference to $y_{\mathrm{Sol}}$.
\end{proof}

%---------------------------------------------------------------------
\section{Assembly}\label{sec:assembly}
%---------------------------------------------------------------------

Theorem~\ref{thm:main} now follows.  Pick a rational $k > k_\star$
(possible by density of $\Q$ in $\R$);
Lemma~\ref{lem:nonneg} gives clause~(i),
Lemma~\ref{lem:sigma} gives clause~(ii),
and Lemma~\ref{lem:identity} gives clause~(iii).

\begin{remark}[Boundary case $n = 0$]
The $n = 0$ case is vacuous: with no variables, all universal
statements are trivially satisfied.
\end{remark}

%---------------------------------------------------------------------
\section{Comparison with the scalar cubic}
%---------------------------------------------------------------------

For $p(y) = 1 - y^3$ on $[-1, 1]$, the threshold
$k_\star = 6$ matches the saddle-node value derived in the companion
scalar-cubic note via direct phase-plane analysis.  The general
construction here gives the same constant (up to the rounding implicit
in $\Q$-density), but the proof is substantially more abstract: it
does not exploit the explicit $1$-dim structure, instead leveraging
only the $\sigma$-Nagumo argument.  The advantage is generality; the
cost is a slightly weaker statement (the box width $2\beta$ vs the
sharper scalar-cubic invariant).

%---------------------------------------------------------------------
\section*{Notes}
%---------------------------------------------------------------------

The full Lean~4 formalization, organized along the lemma sequence
above (algebraic identities, threshold witness, Picard existence, the
three-pronged invariant, dual-rail identity, main theorem), is available
in the Ripple repository under
\path{Ripple/DualRail/ConstantAnnihilationGeneral.lean}.  The proof
depends only on Lean's standard axiomatic core (propositional
extensionality, choice, soundness of quotients).

The companion note on the scalar cubic case ($p(y) = 1 - y^3$,
threshold $k_\star = 6$) is at
\path{notes/scalar-cubic-dual-rail.tex}.

\end{document}
